%%%%%%%%%%%%%%%%%%%%%%%%%%%%%%%%%%%%%%%%%%%%%%%%%%%%%%%%%%%%
% 1.15 LATTICE THEORY
% The Composition of Advanced Mathematics
%%%%%%%%%%%%%%%%%%%%%%%%%%%%%%%%%%%%%%%%%%%%%%%%%%%%%%%%%%%%

\section{Lattice Theory}
\label{sec:lattice-theory}

\prerequisite{
Relations and partial orders from \cref{sec:relations}, elementary set theory
from \cref{sec:elementary-set-theory}, abstract algebra from
\cref{sec:abstract-algebra}, and universal algebra from
\cref{sec:universal-algebra}.
}

\learningobjectives{
By the end of this section, the reader should be able to:
\begin{itemize}
    \item distinguish partially ordered sets from lattices;
    \item identify comparable and incomparable elements;
    \item construct and interpret Hasse diagrams;
    \item determine minimal, maximal, least, and greatest elements;
    \item determine upper bounds, lower bounds, suprema, and infima;
    \item define lattices using both order-theoretic and algebraic viewpoints;
    \item compute meets and joins;
    \item recognize chains, antichains, sublattices, and product lattices;
    \item distinguish bounded, complemented, distributive, and modular lattices;
    \item work with complete lattices;
    \item understand Boolean algebras as complemented distributive lattices;
    \item define lattice homomorphisms and lattice isomorphisms;
    \item interpret congruence lattices;
    \item state the Knaster--Tarski fixed-point theorem;
    \item explain connections with algebra, logic, topology, and computer science.
\end{itemize}
}

%%%%%%%%%%%%%%%%%%%%%%%%%%%%%%%%%%%%%%%%%%%%%%%%%%%%%%%%%%%%
\subsection{From Order to Lattice}
\label{subsec:from-order-to-lattice}
%%%%%%%%%%%%%%%%%%%%%%%%%%%%%%%%%%%%%%%%%%%%%%%%%%%%%%%%%%%%

A partially ordered set provides a way to compare elements without requiring
every pair to be comparable.

Lattice theory begins with this order structure and asks a stronger question:

\begin{quote}
Given two elements, is there a best element above both and a best element
below both?
\end{quote}

When both always exist, the partially ordered set is a lattice.

\begin{conceptbox}
A \textbf{lattice} is an ordered structure in which every pair of elements
has both a least upper bound and a greatest lower bound.
\end{conceptbox}

Lattices appear throughout mathematics because many different notions of
``combining upward'' and ``combining downward'' satisfy this pattern.

Examples include:

\begin{itemize}
    \item union and intersection of sets;
    \item greatest common divisors and least common multiples;
    \item logical OR and AND;
    \item subgroups ordered by inclusion;
    \item subspaces ordered by inclusion;
    \item divisibility among integers;
    \item congruence relations on an algebra.
\end{itemize}

%%%%%%%%%%%%%%%%%%%%%%%%%%%%%%%%%%%%%%%%%%%%%%%%%%%%%%%%%%%%
\subsection{Review of Partial Orders}
\label{subsec:lattice-partial-orders}
%%%%%%%%%%%%%%%%%%%%%%%%%%%%%%%%%%%%%%%%%%%%%%%%%%%%%%%%%%%%

A relation \(\leq\) on a set \(P\) is a partial order if it is reflexive,
antisymmetric, and transitive.

Thus:

\[ x\leq x, \]

\[ x\leq y\text{ and }y\leq x\Longrightarrow x=y, \]

and

\[ x\leq y\text{ and }y\leq z\Longrightarrow x\leq z. \]

The pair

\[ (P,\leq) \]

is called a \emph{partially ordered set}, or \emph{poset}.

The notation is read ``\(P\) partially ordered by less than or equal to.''

The symbol \(\leq\) is abstract here. It need not represent ordinary
numerical inequality.

%%%%%%%%%%%%%%%%%%%%%%%%%%%%%%%%%%%%%%%%%%%%%%%%%%%%%%%%%%%%
\subsection{Comparable and Incomparable Elements}
\label{subsec:comparable-incomparable}
%%%%%%%%%%%%%%%%%%%%%%%%%%%%%%%%%%%%%%%%%%%%%%%%%%%%%%%%%%%%

Two elements \(x,y\in P\) are \emph{comparable} if

\[ x\leq y\qquad\text{or}\qquad y\leq x. \]

Otherwise they are \emph{incomparable}.

Incomparability is often denoted

\[ x\parallel y. \]

The symbol \(\parallel\) is read here as ``\(x\) is incomparable with
\(y\).''

For example, under divisibility,

\[ 2\mid6, \]

so \(2\) and \(6\) are comparable.

However,

\[ 2\nmid3,\qquad3\nmid2, \]

so \(2\) and \(3\) are incomparable.

%%%%%%%%%%%%%%%%%%%%%%%%%%%%%%%%%%%%%%%%%%%%%%%%%%%%%%%%%%%%
\subsection{Chains and Antichains}
\label{subsec:chains-antichains}
%%%%%%%%%%%%%%%%%%%%%%%%%%%%%%%%%%%%%%%%%%%%%%%%%%%%%%%%%%%%

\begin{definitionbox}{Chain}
A \emph{chain} in a poset is a subset in which every pair of elements is
comparable.
\end{definitionbox}

Thus a chain behaves like a totally ordered subset.

\begin{definitionbox}{Antichain}
An \emph{antichain} is a subset in which no two distinct elements are
comparable.
\end{definitionbox}

For the power set

\[ \mathcal P(\{1,2,3\}) \]

ordered by inclusion,

\[ \varnothing\subseteq\{1\}\subseteq\{1,2\}\subseteq\{1,2,3\} \]

is a chain, while

\[ \bigl\{\{1\},\{2\},\{3\}\bigr\} \]

is an antichain.

%%%%%%%%%%%%%%%%%%%%%%%%%%%%%%%%%%%%%%%%%%%%%%%%%%%%%%%%%%%%
\subsection{Covering Relations}
\label{subsec:covering-relations}
%%%%%%%%%%%%%%%%%%%%%%%%%%%%%%%%%%%%%%%%%%%%%%%%%%%%%%%%%%%%

\begin{definitionbox}{Cover}
Let \(x<y\) in a poset. We say that \(y\) \emph{covers} \(x\) if there is no
element \(z\) satisfying
\[ x<z<y. \]
\end{definitionbox}

A covering relation records an immediate step upward in the order.

One commonly writes

\[ x\prec y. \]

The symbol \(\prec\) is read ``\(x\) is covered by \(y\)'' in this context.

%%%%%%%%%%%%%%%%%%%%%%%%%%%%%%%%%%%%%%%%%%%%%%%%%%%%%%%%%%%%
\subsection{Hasse Diagrams}
\label{subsec:hasse-diagrams}
%%%%%%%%%%%%%%%%%%%%%%%%%%%%%%%%%%%%%%%%%%%%%%%%%%%%%%%%%%%%

A finite poset can often be represented by a \emph{Hasse diagram}.

In a Hasse diagram:

\begin{itemize}
    \item larger elements are placed above smaller elements;
    \item an edge connects \(x\) to \(y\) when \(y\) covers \(x\);
    \item arrows are usually omitted because upward direction indicates order;
    \item reflexive and transitive edges are omitted.
\end{itemize}

For example, consider the divisors of \(12\),

\[ \{1,2,3,4,6,12\}, \]

ordered by divisibility.

\begin{center}
\begin{tikzpicture}[
    node distance=12mm and 16mm,
    every node/.style={draw,circle,minimum size=8mm,inner sep=1pt}
]
\node (one) {\(1\)};
\node[above left=of one] (two) {\(2\)};
\node[above right=of one] (three) {\(3\)};
\node[above left=of two] (four) {\(4\)};
\node[above right=of two] (six) {\(6\)};
\node[above=of six] (twelve) {\(12\)};

\draw (one)--(two);
\draw (one)--(three);
\draw (two)--(four);
\draw (two)--(six);
\draw (three)--(six);
\draw (four)--(twelve);
\draw (six)--(twelve);
\end{tikzpicture}
\end{center}

The diagram encodes the divisibility order without drawing every implied
relation.

%%%%%%%%%%%%%%%%%%%%%%%%%%%%%%%%%%%%%%%%%%%%%%%%%%%%%%%%%%%%
\subsection{Minimal and Maximal Elements}
\label{subsec:minimal-maximal}
%%%%%%%%%%%%%%%%%%%%%%%%%%%%%%%%%%%%%%%%%%%%%%%%%%%%%%%%%%%%

\begin{definitionbox}{Minimal Element}
An element \(m\in P\) is \emph{minimal} if there is no \(x\in P\) such that
\[ x<m. \]
\end{definitionbox}

\begin{definitionbox}{Maximal Element}
An element \(m\in P\) is \emph{maximal} if there is no \(x\in P\) such that
\[ m<x. \]
\end{definitionbox}

A poset may have several minimal or maximal elements.

This differs from least and greatest elements.

%%%%%%%%%%%%%%%%%%%%%%%%%%%%%%%%%%%%%%%%%%%%%%%%%%%%%%%%%%%%
\subsection{Least and Greatest Elements}
\label{subsec:least-greatest-elements}
%%%%%%%%%%%%%%%%%%%%%%%%%%%%%%%%%%%%%%%%%%%%%%%%%%%%%%%%%%%%

\begin{definitionbox}{Least Element}
An element \(0\in P\) is a \emph{least element} if
\[ 0\leq x \]
for every \(x\in P\).
\end{definitionbox}

\begin{definitionbox}{Greatest Element}
An element \(1\in P\) is a \emph{greatest element} if
\[ x\leq1 \]
for every \(x\in P\).
\end{definitionbox}

The symbols \(0\) and \(1\) are conventional lattice notation and do not
necessarily represent numerical zero and one.

A least element is automatically the unique minimal element.

A greatest element is automatically the unique maximal element.

The converses need not hold.

%%%%%%%%%%%%%%%%%%%%%%%%%%%%%%%%%%%%%%%%%%%%%%%%%%%%%%%%%%%%
\subsection{Minimal Does Not Mean Least}
\label{subsec:minimal-not-least}
%%%%%%%%%%%%%%%%%%%%%%%%%%%%%%%%%%%%%%%%%%%%%%%%%%%%%%%%%%%%

Consider

\begin{center}
\begin{tikzpicture}[
    node distance=14mm and 18mm,
    every node/.style={draw,circle,minimum size=8mm}
]
\node (a) {\(a\)};
\node[right=of a] (b) {\(b\)};
\node[above=of $(a)!0.5!(b)$] (c) {\(c\)};

\draw (a)--(c);
\draw (b)--(c);
\end{tikzpicture}
\end{center}

Both \(a\) and \(b\) are minimal.

Neither is least because neither lies below the other.

The element \(c\) is greatest.

This distinction is fundamental in order theory.

%%%%%%%%%%%%%%%%%%%%%%%%%%%%%%%%%%%%%%%%%%%%%%%%%%%%%%%%%%%%
\subsection{Upper and Lower Bounds}
\label{subsec:upper-lower-bounds}
%%%%%%%%%%%%%%%%%%%%%%%%%%%%%%%%%%%%%%%%%%%%%%%%%%%%%%%%%%%%

Let

\[ S\subseteq P. \]

\begin{definitionbox}{Upper Bound}
An element \(u\in P\) is an \emph{upper bound} of \(S\) if
\[ s\leq u \]
for every \(s\in S\).
\end{definitionbox}

\begin{definitionbox}{Lower Bound}
An element \(\ell\in P\) is a \emph{lower bound} of \(S\) if
\[ \ell\leq s \]
for every \(s\in S\).
\end{definitionbox}

An upper or lower bound need not belong to \(S\).

%%%%%%%%%%%%%%%%%%%%%%%%%%%%%%%%%%%%%%%%%%%%%%%%%%%%%%%%%%%%
\subsection{Suprema and Infima}
\label{subsec:suprema-infima-lattice}
%%%%%%%%%%%%%%%%%%%%%%%%%%%%%%%%%%%%%%%%%%%%%%%%%%%%%%%%%%%%

\begin{definitionbox}{Supremum}
The \emph{supremum} of \(S\), if it exists, is the least upper bound of
\(S\).
\end{definitionbox}

It is commonly written

\[ \sup S. \]

This is read ``the supremum of \(S\)'' or ``the least upper bound of \(S\).''

\begin{definitionbox}{Infimum}
The \emph{infimum} of \(S\), if it exists, is the greatest lower bound of
\(S\).
\end{definitionbox}

It is written

\[ \inf S. \]

This is read ``the infimum of \(S\)'' or ``the greatest lower bound of
\(S\).''

These notions later play a major role in real analysis.

%%%%%%%%%%%%%%%%%%%%%%%%%%%%%%%%%%%%%%%%%%%%%%%%%%%%%%%%%%%%
\subsection{Join and Meet}
\label{subsec:join-meet}
%%%%%%%%%%%%%%%%%%%%%%%%%%%%%%%%%%%%%%%%%%%%%%%%%%%%%%%%%%%%

For two elements \(x,y\), lattice theory uses special notation for their
least upper bound and greatest lower bound.

\begin{definitionbox}{Join}
The \emph{join} of \(x\) and \(y\) is their least upper bound:
\[ x\vee y=\sup\{x,y\}. \]
\end{definitionbox}

The symbol \(\vee\) is read ``join.''

Thus

\[ x\vee y \]

is read ``\(x\) join \(y\).''

\begin{definitionbox}{Meet}
The \emph{meet} of \(x\) and \(y\) is their greatest lower bound:
\[ x\wedge y=\inf\{x,y\}. \]
\end{definitionbox}

The symbol \(\wedge\) is read ``meet.''

Thus

\[ x\wedge y \]

is read ``\(x\) meet \(y\).''

%%%%%%%%%%%%%%%%%%%%%%%%%%%%%%%%%%%%%%%%%%%%%%%%%%%%%%%%%%%%
\subsection{Visual Meaning of Meet and Join}
\label{subsec:visual-meet-join}
%%%%%%%%%%%%%%%%%%%%%%%%%%%%%%%%%%%%%%%%%%%%%%%%%%%%%%%%%%%%

Suppose a lattice contains elements arranged as

\begin{center}
\begin{tikzpicture}[
    node distance=12mm and 18mm,
    every node/.style={draw,circle,minimum size=10mm}
]
\node (meet) {\(x\wedge y\)};
\node[above left=of meet] (x) {\(x\)};
\node[above right=of meet] (y) {\(y\)};
\node[above=22mm of meet] (join) {\(x\vee y\)};

\draw (meet)--(x);
\draw (meet)--(y);
\draw (x)--(join);
\draw (y)--(join);
\end{tikzpicture}
\end{center}

The meet is the highest element lying below both \(x\) and \(y\).

The join is the lowest element lying above both \(x\) and \(y\).

%%%%%%%%%%%%%%%%%%%%%%%%%%%%%%%%%%%%%%%%%%%%%%%%%%%%%%%%%%%%
\subsection{Order-Theoretic Definition of a Lattice}
\label{subsec:order-lattice-definition}
%%%%%%%%%%%%%%%%%%%%%%%%%%%%%%%%%%%%%%%%%%%%%%%%%%%%%%%%%%%%

\begin{definitionbox}{Lattice}
A poset \((L,\leq)\) is a \emph{lattice} if every pair \(x,y\in L\) has a
meet and a join.
\end{definitionbox}

Equivalently, for every \(x,y\in L\),

\[ x\wedge y\quad\text{and}\quad x\vee y \]

must exist.

The requirement concerns pairs of elements, not necessarily arbitrary
subsets.

%%%%%%%%%%%%%%%%%%%%%%%%%%%%%%%%%%%%%%%%%%%%%%%%%%%%%%%%%%%%
\subsection{A Poset That Is Not a Lattice}
\label{subsec:poset-not-lattice}
%%%%%%%%%%%%%%%%%%%%%%%%%%%%%%%%%%%%%%%%%%%%%%%%%%%%%%%%%%%%

Consider

\begin{center}
\begin{tikzpicture}[
    node distance=14mm and 20mm,
    every node/.style={draw,circle,minimum size=8mm}
]
\node (a) {\(a\)};
\node[right=of a] (b) {\(b\)};
\node[above=of a] (c) {\(c\)};
\node[above=of b] (d) {\(d\)};

\draw (a)--(c);
\draw (b)--(d);
\end{tikzpicture}
\end{center}

The elements \(a\) and \(b\) have no common upper bound.

Therefore \(a\vee b\) does not exist.

The poset is not a lattice.

%%%%%%%%%%%%%%%%%%%%%%%%%%%%%%%%%%%%%%%%%%%%%%%%%%%%%%%%%%%%
\subsection{The Power-Set Lattice}
\label{subsec:power-set-lattice}
%%%%%%%%%%%%%%%%%%%%%%%%%%%%%%%%%%%%%%%%%%%%%%%%%%%%%%%%%%%%

Let \(X\) be a set.

The power set

\[ \mathcal P(X) \]

ordered by inclusion is a lattice.

For subsets \(A,B\subseteq X\),

\[ A\wedge B=A\cap B, \qquad A\vee B=A\cup B. \]

Thus intersection acts as meet and union acts as join.

For

\[ X=\{1,2,3\}, \]

the lattice has the Hasse diagram

\begin{center}
\begin{tikzpicture}[
    node distance=11mm and 13mm,
    every node/.style={draw,rounded corners,inner sep=3pt,font=\small}
]
\node (empty) {\(\varnothing\)};

\node[above left=of empty] (one) {\(\{1\}\)};
\node[above=of empty] (two) {\(\{2\}\)};
\node[above right=of empty] (three) {\(\{3\}\)};

\node[above left=15mm and -2mm of two] (onetwo) {\(\{1,2\}\)};
\node[above=15mm of two] (onethree) {\(\{1,3\}\)};
\node[above right=15mm and -2mm of two] (twothree) {\(\{2,3\}\)};

\node[above=15mm of onethree] (all) {\(\{1,2,3\}\)};

\draw (empty)--(one);
\draw (empty)--(two);
\draw (empty)--(three);

\draw (one)--(onetwo);
\draw (one)--(onethree);

\draw (two)--(onetwo);
\draw (two)--(twothree);

\draw (three)--(onethree);
\draw (three)--(twothree);

\draw (onetwo)--(all);
\draw (onethree)--(all);
\draw (twothree)--(all);
\end{tikzpicture}
\end{center}

This lattice is often denoted

\[ B_3. \]

%%%%%%%%%%%%%%%%%%%%%%%%%%%%%%%%%%%%%%%%%%%%%%%%%%%%%%%%%%%%
\subsection{Divisibility Lattices}
\label{subsec:divisibility-lattices}
%%%%%%%%%%%%%%%%%%%%%%%%%%%%%%%%%%%%%%%%%%%%%%%%%%%%%%%%%%%%

Let

\[ D_n=\{d\in\N:d\mid n\} \]

be the set of positive divisors of \(n\), ordered by divisibility.

Then

\[ a\wedge b=\gcd(a,b), \]

and

\[ a\vee b=\operatorname{lcm}(a,b). \]

Thus greatest common divisor and least common multiple are lattice
operations.

For example, among divisors of \(60\),

\[ 12\wedge20=\gcd(12,20)=4, \]

and

\[ 12\vee20=\operatorname{lcm}(12,20)=60. \]

%%%%%%%%%%%%%%%%%%%%%%%%%%%%%%%%%%%%%%%%%%%%%%%%%%%%%%%%%%%%
\subsection{Worked Example: Meet and Join}
\label{subsec:worked-meet-join}
%%%%%%%%%%%%%%%%%%%%%%%%%%%%%%%%%%%%%%%%%%%%%%%%%%%%%%%%%%%%

\begin{workedexamplebox}{Meet and Join in a Divisibility Lattice}

Consider the divisors of \(36\), ordered by divisibility. Determine the meet
and join of \(6\) and \(9\).

\begin{tabularx}{\textwidth}{@{}>{\raggedright\arraybackslash}p{0.42\textwidth}X@{}}
The meet is the greatest common divisor. &
\(6\wedge9=\gcd(6,9)=3\)\\
The join is the least common multiple. &
\(6\vee9=\operatorname{lcm}(6,9)=18\)
\end{tabularx}

\begin{answerbox}
\[ \boxed{6\wedge9=3,\qquad6\vee9=18.} \]
\end{answerbox}
\end{workedexamplebox}

%%%%%%%%%%%%%%%%%%%%%%%%%%%%%%%%%%%%%%%%%%%%%%%%%%%%%%%%%%%%
\subsection{Subspace Lattices}
\label{subsec:subspace-lattices}
%%%%%%%%%%%%%%%%%%%%%%%%%%%%%%%%%%%%%%%%%%%%%%%%%%%%%%%%%%%%

Let \(V\) be a vector space.

The collection of subspaces of \(V\), ordered by inclusion, forms a lattice.

For subspaces \(U,W\leq V\),

\[ U\wedge W=U\cap W, \]

while

\[ U\vee W=U+W. \]

The sum \(U+W\) is the smallest subspace containing both \(U\) and \(W\).

This gives an important connection between lattice theory and linear
algebra.

%%%%%%%%%%%%%%%%%%%%%%%%%%%%%%%%%%%%%%%%%%%%%%%%%%%%%%%%%%%%
\subsection{Subgroup Lattices}
\label{subsec:subgroup-lattices}
%%%%%%%%%%%%%%%%%%%%%%%%%%%%%%%%%%%%%%%%%%%%%%%%%%%%%%%%%%%%

For a group \(G\), the set of subgroups

\[ \operatorname{Sub}(G) \]

ordered by inclusion forms a lattice.

For subgroups \(H,K\leq G\),

\[ H\wedge K=H\cap K. \]

The join is the subgroup generated by both:

\[ H\vee K=\langle H\cup K\rangle. \]

It is important that \(H\cup K\) itself is generally not a subgroup.

%%%%%%%%%%%%%%%%%%%%%%%%%%%%%%%%%%%%%%%%%%%%%%%%%%%%%%%%%%%%
\subsection{The Algebraic Definition of a Lattice}
\label{subsec:algebraic-lattice-definition}
%%%%%%%%%%%%%%%%%%%%%%%%%%%%%%%%%%%%%%%%%%%%%%%%%%%%%%%%%%%%

A lattice can also be defined without explicitly mentioning an order.

\begin{definitionbox}{Lattice Algebra}
A \emph{lattice algebra} is a set \(L\) equipped with two binary operations
\(\wedge\) and \(\vee\) satisfying the commutative, associative,
idempotent, and absorption identities.
\end{definitionbox}

These identities are

\[
x\wedge y=y\wedge x,\qquad
x\vee y=y\vee x,
\]

\[
(x\wedge y)\wedge z=x\wedge(y\wedge z),\qquad
(x\vee y)\vee z=x\vee(y\vee z),
\]

\[
x\wedge x=x,\qquad x\vee x=x,
\]

and

\[
x\wedge(x\vee y)=x,\qquad
x\vee(x\wedge y)=x.
\]

%%%%%%%%%%%%%%%%%%%%%%%%%%%%%%%%%%%%%%%%%%%%%%%%%%%%%%%%%%%%
\subsection{Absorption Laws}
\label{subsec:absorption-laws}
%%%%%%%%%%%%%%%%%%%%%%%%%%%%%%%%%%%%%%%%%%%%%%%%%%%%%%%%%%%%

The identities

\[ x\wedge(x\vee y)=x \]

and

\[ x\vee(x\wedge y)=x \]

are called the \emph{absorption laws}.

They encode the relationship between meet and join.

If \(x\vee y\) already lies above \(x\), then meeting it back with \(x\)
returns \(x\).

Dually, if \(x\wedge y\) lies below \(x\), joining it with \(x\) returns
\(x\).

%%%%%%%%%%%%%%%%%%%%%%%%%%%%%%%%%%%%%%%%%%%%%%%%%%%%%%%%%%%%
\subsection{Order from Meet and Join}
\label{subsec:order-from-meet-join}
%%%%%%%%%%%%%%%%%%%%%%%%%%%%%%%%%%%%%%%%%%%%%%%%%%%%%%%%%%%%

Given a lattice algebra, define

\[ x\leq y\iff x\wedge y=x. \]

Equivalently,

\[ x\leq y\iff x\vee y=y. \]

These conditions recover the underlying partial order.

\begin{theorem}
For a lattice,
\[ x\leq y
\iff x\wedge y=x
\iff x\vee y=y. \]
\end{theorem}

This establishes the equivalence between the order-theoretic and algebraic
definitions of a lattice.

%%%%%%%%%%%%%%%%%%%%%%%%%%%%%%%%%%%%%%%%%%%%%%%%%%%%%%%%%%%%
\subsection{Duality Principle}
\label{subsec:lattice-duality}
%%%%%%%%%%%%%%%%%%%%%%%%%%%%%%%%%%%%%%%%%%%%%%%%%%%%%%%%%%%%

Lattice theory possesses a powerful symmetry.

If every order relation is reversed, then meets and joins exchange roles.

Thus

\[ \leq\longleftrightarrow\geq, \qquad
\wedge\longleftrightarrow\vee, \qquad
0\longleftrightarrow1. \]

\begin{principle}[Principle of Duality]
If a statement about lattices follows from the lattice axioms, then the
dual statement obtained by interchanging meet and join and reversing the
order also holds.
\end{principle}

For example,

\[ x\wedge(x\vee y)=x \]

has dual

\[ x\vee(x\wedge y)=x. \]

%%%%%%%%%%%%%%%%%%%%%%%%%%%%%%%%%%%%%%%%%%%%%%%%%%%%%%%%%%%%
\subsection{Sublattices}
\label{subsec:sublattices}
%%%%%%%%%%%%%%%%%%%%%%%%%%%%%%%%%%%%%%%%%%%%%%%%%%%%%%%%%%%%

\begin{definitionbox}{Sublattice}
A nonempty subset \(S\subseteq L\) is a \emph{sublattice} if it is closed
under meet and join.
\end{definitionbox}

Thus for every \(x,y\in S\),

\[ x\wedge y\in S,\qquad x\vee y\in S. \]

A subset can be a subposet without being a sublattice.

The additional requirement is closure under the lattice operations.

%%%%%%%%%%%%%%%%%%%%%%%%%%%%%%%%%%%%%%%%%%%%%%%%%%%%%%%%%%%%
\subsection{Generated Sublattices}
\label{subsec:generated-sublattices}
%%%%%%%%%%%%%%%%%%%%%%%%%%%%%%%%%%%%%%%%%%%%%%%%%%%%%%%%%%%%

For

\[ X\subseteq L, \]

the smallest sublattice containing \(X\) is denoted

\[ \langle X\rangle_{\mathrm{Lat}}. \]

It is obtained by repeatedly applying meet and join to elements of \(X\).

This is a special case of the generated-subalgebra construction from
universal algebra.

%%%%%%%%%%%%%%%%%%%%%%%%%%%%%%%%%%%%%%%%%%%%%%%%%%%%%%%%%%%%
\subsection{Product Lattices}
\label{subsec:product-lattices}
%%%%%%%%%%%%%%%%%%%%%%%%%%%%%%%%%%%%%%%%%%%%%%%%%%%%%%%%%%%%

If \(L\) and \(M\) are lattices, their Cartesian product

\[ L\times M \]

becomes a lattice under coordinatewise order:

\[ (x_1,y_1)\leq(x_2,y_2)
\iff x_1\leq x_2\text{ and }y_1\leq y_2. \]

Meet and join are computed coordinatewise:

\[ (x_1,y_1)\wedge(x_2,y_2)
=(x_1\wedge x_2,\ y_1\wedge y_2), \]

\[ (x_1,y_1)\vee(x_2,y_2)
=(x_1\vee x_2,\ y_1\vee y_2). \]

%%%%%%%%%%%%%%%%%%%%%%%%%%%%%%%%%%%%%%%%%%%%%%%%%%%%%%%%%%%%
\subsection{Bounded Lattices}
\label{subsec:bounded-lattices}
%%%%%%%%%%%%%%%%%%%%%%%%%%%%%%%%%%%%%%%%%%%%%%%%%%%%%%%%%%%%

\begin{definitionbox}{Bounded Lattice}
A lattice \(L\) is \emph{bounded} if it contains both a least element \(0\)
and a greatest element \(1\).
\end{definitionbox}

Thus

\[ 0\leq x\leq1 \]

for every \(x\in L\).

Equivalently,

\[ 0\wedge x=0,\qquad0\vee x=x, \]

and

\[ 1\wedge x=x,\qquad1\vee x=1. \]

%%%%%%%%%%%%%%%%%%%%%%%%%%%%%%%%%%%%%%%%%%%%%%%%%%%%%%%%%%%%
\subsection{Complements}
\label{subsec:lattice-complements}
%%%%%%%%%%%%%%%%%%%%%%%%%%%%%%%%%%%%%%%%%%%%%%%%%%%%%%%%%%%%

\begin{definitionbox}{Complement}
Let \(L\) be a bounded lattice. An element \(y\) is a \emph{complement} of
\(x\) if
\[ x\wedge y=0,\qquad x\vee y=1. \]
\end{definitionbox}

A complement is often written

\[ x'. \]

This is read ``\(x\) complement.''

Depending on the lattice, an element may have:

\begin{itemize}
    \item no complement;
    \item exactly one complement;
    \item several complements.
\end{itemize}

%%%%%%%%%%%%%%%%%%%%%%%%%%%%%%%%%%%%%%%%%%%%%%%%%%%%%%%%%%%%
\subsection{Complemented Lattices}
\label{subsec:complemented-lattices}
%%%%%%%%%%%%%%%%%%%%%%%%%%%%%%%%%%%%%%%%%%%%%%%%%%%%%%%%%%%%

\begin{definitionbox}{Complemented Lattice}
A bounded lattice is \emph{complemented} if every element has at least one
complement.
\end{definitionbox}

In the power-set lattice,

\[ \mathcal P(X), \]

the complement of \(A\subseteq X\) is

\[ A^c=X\setminus A. \]

Indeed,

\[ A\cap A^c=\varnothing, \qquad A\cup A^c=X. \]

%%%%%%%%%%%%%%%%%%%%%%%%%%%%%%%%%%%%%%%%%%%%%%%%%%%%%%%%%%%%
\subsection{Distributive Lattices}
\label{subsec:distributive-lattices}
%%%%%%%%%%%%%%%%%%%%%%%%%%%%%%%%%%%%%%%%%%%%%%%%%%%%%%%%%%%%

\begin{definitionbox}{Distributive Lattice}
A lattice is \emph{distributive} if
\[ x\wedge(y\vee z)
=(x\wedge y)\vee(x\wedge z) \]
and
\[ x\vee(y\wedge z)
=(x\vee y)\wedge(x\vee z) \]
for all \(x,y,z\in L\).
\end{definitionbox}

By lattice duality, either distributive identity implies the other in a
lattice.

The power-set lattice is distributive because set intersection and union
satisfy the familiar distributive laws.

%%%%%%%%%%%%%%%%%%%%%%%%%%%%%%%%%%%%%%%%%%%%%%%%%%%%%%%%%%%%
\subsection{The Diamond Lattice \(M_3\)}
\label{subsec:diamond-lattice}
%%%%%%%%%%%%%%%%%%%%%%%%%%%%%%%%%%%%%%%%%%%%%%%%%%%%%%%%%%%%

A standard nondistributive lattice is the five-element diamond lattice
\(M_3\).

\begin{center}
\begin{tikzpicture}[
    node distance=13mm and 15mm,
    every node/.style={draw,circle,minimum size=8mm}
]
\node (zero) {\(0\)};
\node[above left=of zero] (a) {\(a\)};
\node[above=of zero] (b) {\(b\)};
\node[above right=of zero] (c) {\(c\)};
\node[above=26mm of zero] (one) {\(1\)};

\draw (zero)--(a);
\draw (zero)--(b);
\draw (zero)--(c);
\draw (a)--(one);
\draw (b)--(one);
\draw (c)--(one);
\end{tikzpicture}
\end{center}

For distinct \(a,b,c\),

\[ b\vee c=1. \]

Therefore

\[ a\wedge(b\vee c)=a. \]

But

\[ a\wedge b=0,\qquad a\wedge c=0, \]

so

\[ (a\wedge b)\vee(a\wedge c)=0. \]

Hence distributivity fails.

%%%%%%%%%%%%%%%%%%%%%%%%%%%%%%%%%%%%%%%%%%%%%%%%%%%%%%%%%%%%
\subsection{The Pentagon Lattice \(N_5\)}
\label{subsec:pentagon-lattice}
%%%%%%%%%%%%%%%%%%%%%%%%%%%%%%%%%%%%%%%%%%%%%%%%%%%%%%%%%%%%

Another standard nondistributive lattice is \(N_5\).

\begin{center}
\begin{tikzpicture}[
    node distance=13mm and 18mm,
    every node/.style={draw,circle,minimum size=8mm}
]
\node (zero) {\(0\)};
\node[above left=of zero] (a) {\(a\)};
\node[above=of a] (b) {\(b\)};
\node[above right=of zero] (c) {\(c\)};
\node[above right=of b] (one) {\(1\)};

\draw (zero)--(a);
\draw (a)--(b);
\draw (b)--(one);
\draw (zero)--(c);
\draw (c)--(one);
\end{tikzpicture}
\end{center}

The shape resembles a pentagon, which motivates its name.

The lattices \(M_3\) and \(N_5\) characterize failure of distributivity in
an important theorem.

%%%%%%%%%%%%%%%%%%%%%%%%%%%%%%%%%%%%%%%%%%%%%%%%%%%%%%%%%%%%
\subsection{Characterization of Distributive Lattices}
\label{subsec:distributive-characterization}
%%%%%%%%%%%%%%%%%%%%%%%%%%%%%%%%%%%%%%%%%%%%%%%%%%%%%%%%%%%%

\begin{theorem}[Birkhoff's Characterization]
A lattice is distributive if and only if it contains no sublattice
isomorphic to \(M_3\) or \(N_5\).
\end{theorem}

Thus the two small lattices \(M_3\) and \(N_5\) are the fundamental
obstructions to distributivity.

%%%%%%%%%%%%%%%%%%%%%%%%%%%%%%%%%%%%%%%%%%%%%%%%%%%%%%%%%%%%
\subsection{Modular Lattices}
\label{subsec:modular-lattices}
%%%%%%%%%%%%%%%%%%%%%%%%%%%%%%%%%%%%%%%%%%%%%%%%%%%%%%%%%%%%

Distributivity is a strong condition. A weaker condition is modularity.

\begin{definitionbox}{Modular Lattice}
A lattice \(L\) is \emph{modular} if whenever \(x\leq z\),
\[ x\vee(y\wedge z)=(x\vee y)\wedge z. \]
\end{definitionbox}

The condition \(x\leq z\) is essential.

Every distributive lattice is modular.

The converse is false.

%%%%%%%%%%%%%%%%%%%%%%%%%%%%%%%%%%%%%%%%%%%%%%%%%%%%%%%%%%%%
\subsection{Distributive Versus Modular}
\label{subsec:distributive-versus-modular}
%%%%%%%%%%%%%%%%%%%%%%%%%%%%%%%%%%%%%%%%%%%%%%%%%%%%%%%%%%%%

The relationship is

\begin{center}
\begin{tikzpicture}[
    node distance=12mm,
    box/.style={draw,rounded corners,align=center,inner sep=6pt},
    >=stealth
]
\node[box] (lat) {Lattices};
\node[box,below=of lat] (mod) {Modular lattices};
\node[box,below=of mod] (dist) {Distributive lattices};

\draw[->] (dist)--node[right]{implies}(mod);
\draw[->] (mod)--node[right]{special case of}(lat);
\end{tikzpicture}
\end{center}

The diamond \(M_3\) is modular but not distributive.

The pentagon \(N_5\) is not modular.

%%%%%%%%%%%%%%%%%%%%%%%%%%%%%%%%%%%%%%%%%%%%%%%%%%%%%%%%%%%%
\subsection{Subspace Lattices Are Modular}
\label{subsec:subspace-modular}
%%%%%%%%%%%%%%%%%%%%%%%%%%%%%%%%%%%%%%%%%%%%%%%%%%%%%%%%%%%%

The lattice of subspaces of a vector space is generally modular.

If \(X\subseteq Z\), then

\[ X+(Y\cap Z)=(X+Y)\cap Z. \]

In lattice notation,

\[ X\vee(Y\wedge Z)=(X\vee Y)\wedge Z. \]

Thus linear algebra provides one of the most important examples of modular
lattices.

The subspace lattice need not be distributive.

%%%%%%%%%%%%%%%%%%%%%%%%%%%%%%%%%%%%%%%%%%%%%%%%%%%%%%%%%%%%
\subsection{Modular Characterization}
\label{subsec:modular-characterization}
%%%%%%%%%%%%%%%%%%%%%%%%%%%%%%%%%%%%%%%%%%%%%%%%%%%%%%%%%%%%

\begin{theorem}
A lattice is modular if and only if it contains no sublattice isomorphic to
\(N_5\).
\end{theorem}

Compare this with distributivity:

\[ \text{distributive}
\iff\text{contains neither }M_3\text{ nor }N_5, \]

while

\[ \text{modular}
\iff\text{contains no }N_5. \]

%%%%%%%%%%%%%%%%%%%%%%%%%%%%%%%%%%%%%%%%%%%%%%%%%%%%%%%%%%%%
\subsection{Boolean Algebras}
\label{subsec:lattice-boolean-algebras}
%%%%%%%%%%%%%%%%%%%%%%%%%%%%%%%%%%%%%%%%%%%%%%%%%%%%%%%%%%%%

\begin{definitionbox}{Boolean Algebra}
A \emph{Boolean algebra} is a bounded distributive lattice in which every
element has a complement.
\end{definitionbox}

Thus Boolean algebra combines:

\begin{itemize}
    \item meet;
    \item join;
    \item least and greatest elements;
    \item distributivity;
    \item complementation.
\end{itemize}

The standard example is

\[ \bigl(\mathcal P(X),\cap,\cup,{}^c,\varnothing,X\bigr). \]

%%%%%%%%%%%%%%%%%%%%%%%%%%%%%%%%%%%%%%%%%%%%%%%%%%%%%%%%%%%%
\subsection{Boolean Operations}
\label{subsec:boolean-operations}
%%%%%%%%%%%%%%%%%%%%%%%%%%%%%%%%%%%%%%%%%%%%%%%%%%%%%%%%%%%%

The same Boolean structure appears in several settings.

\begin{center}
\begin{tabularx}{0.92\textwidth}{@{}lXXX@{}}
\toprule
\textbf{Lattice} & \textbf{Sets} & \textbf{Logic} & \textbf{Bits}\\
\midrule
\(x\wedge y\) & \(A\cap B\) & AND & \(xy\)\\
\(x\vee y\) & \(A\cup B\) & OR & max-like OR\\
\(x'\) & \(A^c\) & NOT & \(1-x\)\\
\(0\) & \(\varnothing\) & false & \(0\)\\
\(1\) & \(X\) & true & \(1\)\\
\bottomrule
\end{tabularx}
\end{center}

This common structure explains the importance of Boolean algebra in logic,
set theory, digital circuits, and computer science.

%%%%%%%%%%%%%%%%%%%%%%%%%%%%%%%%%%%%%%%%%%%%%%%%%%%%%%%%%%%%
\subsection{De Morgan's Laws}
\label{subsec:lattice-demorgan}
%%%%%%%%%%%%%%%%%%%%%%%%%%%%%%%%%%%%%%%%%%%%%%%%%%%%%%%%%%%%

In a Boolean algebra,

\[ (x\wedge y)'=x'\vee y', \]

and

\[ (x\vee y)'=x'\wedge y'. \]

These are \emph{De Morgan's laws}.

They generalize the familiar set identities

\[ (A\cap B)^c=A^c\cup B^c, \]

\[ (A\cup B)^c=A^c\cap B^c. \]

%%%%%%%%%%%%%%%%%%%%%%%%%%%%%%%%%%%%%%%%%%%%%%%%%%%%%%%%%%%%
\subsection{Uniqueness of Complements}
\label{subsec:unique-complements}
%%%%%%%%%%%%%%%%%%%%%%%%%%%%%%%%%%%%%%%%%%%%%%%%%%%%%%%%%%%%

In an arbitrary complemented lattice, complements need not be unique.

In a distributive complemented lattice, however, they are unique.

\begin{theorem}
Every element of a Boolean algebra has a unique complement.
\end{theorem}

\begin{proof}
Suppose \(y\) and \(z\) are complements of \(x\). Then

\[ x\wedge y=0,\qquad x\vee y=1, \]

and

\[ x\wedge z=0,\qquad x\vee z=1. \]

Using distributivity,

\[
\begin{aligned}
y
&=y\wedge1\\
&=y\wedge(x\vee z)\\
&=(y\wedge x)\vee(y\wedge z)\\
&=0\vee(y\wedge z)\\
&=y\wedge z.
\end{aligned}
\]

Similarly,

\[ z=y\wedge z. \]

Therefore \(y=z\).
\end{proof}

%%%%%%%%%%%%%%%%%%%%%%%%%%%%%%%%%%%%%%%%%%%%%%%%%%%%%%%%%%%%
\subsection{Complete Lattices}
\label{subsec:complete-lattices}
%%%%%%%%%%%%%%%%%%%%%%%%%%%%%%%%%%%%%%%%%%%%%%%%%%%%%%%%%%%%

A lattice guarantees meets and joins for pairs.

A complete lattice requires them for arbitrary subsets.

\begin{definitionbox}{Complete Lattice}
A lattice \(L\) is \emph{complete} if every subset \(S\subseteq L\) has both
a supremum and an infimum in \(L\).
\end{definitionbox}

Thus

\[ \bigvee S \]

denotes the join or supremum of all elements of \(S\), while

\[ \bigwedge S \]

denotes their meet or infimum.

These are read ``the join over \(S\)'' and ``the meet over \(S\).''

%%%%%%%%%%%%%%%%%%%%%%%%%%%%%%%%%%%%%%%%%%%%%%%%%%%%%%%%%%%%
\subsection{Empty Meets and Joins}
\label{subsec:empty-meets-joins}
%%%%%%%%%%%%%%%%%%%%%%%%%%%%%%%%%%%%%%%%%%%%%%%%%%%%%%%%%%%%

In a complete lattice, even the empty subset must have a supremum and an
infimum.

The conventions are

\[ \bigvee\varnothing=0, \qquad \bigwedge\varnothing=1. \]

Therefore every complete lattice is bounded.

%%%%%%%%%%%%%%%%%%%%%%%%%%%%%%%%%%%%%%%%%%%%%%%%%%%%%%%%%%%%
\subsection{Power Sets Are Complete Lattices}
\label{subsec:powerset-complete}
%%%%%%%%%%%%%%%%%%%%%%%%%%%%%%%%%%%%%%%%%%%%%%%%%%%%%%%%%%%%

For any set \(X\),

\[ \mathcal P(X) \]

is a complete lattice under inclusion.

For any family

\[ \mathcal A\subseteq\mathcal P(X), \]

we have

\[ \bigvee\mathcal A=\bigcup_{A\in\mathcal A}A, \]

and

\[ \bigwedge\mathcal A=\bigcap_{A\in\mathcal A}A. \]

Thus arbitrary unions and intersections provide arbitrary joins and meets.

%%%%%%%%%%%%%%%%%%%%%%%%%%%%%%%%%%%%%%%%%%%%%%%%%%%%%%%%%%%%
\subsection{Finite Lattices and Completeness}
\label{subsec:finite-lattice-completeness}
%%%%%%%%%%%%%%%%%%%%%%%%%%%%%%%%%%%%%%%%%%%%%%%%%%%%%%%%%%%%

A finite lattice is complete if it is bounded.

In particular, every finite nonempty lattice has joins and meets of every
nonempty finite subset.

If it also possesses \(0\) and \(1\), then the empty join and meet exist as
well.

%%%%%%%%%%%%%%%%%%%%%%%%%%%%%%%%%%%%%%%%%%%%%%%%%%%%%%%%%%%%
\subsection{Lattice Homomorphisms}
\label{subsec:lattice-homomorphisms}
%%%%%%%%%%%%%%%%%%%%%%%%%%%%%%%%%%%%%%%%%%%%%%%%%%%%%%%%%%%%

\begin{definitionbox}{Lattice Homomorphism}
Let \(L\) and \(M\) be lattices. A function
\[ f:L\to M \]
is a \emph{lattice homomorphism} if
\[ f(x\wedge y)=f(x)\wedge f(y) \]
and
\[ f(x\vee y)=f(x)\vee f(y) \]
for all \(x,y\in L\).
\end{definitionbox}

Thus a lattice homomorphism preserves both fundamental lattice operations.

For bounded lattices, one may additionally require

\[ f(0)=0,\qquad f(1)=1, \]

depending on the chosen signature.

%%%%%%%%%%%%%%%%%%%%%%%%%%%%%%%%%%%%%%%%%%%%%%%%%%%%%%%%%%%%
\subsection{Lattice Homomorphisms Preserve Order}
\label{subsec:lattice-hom-preserve-order}
%%%%%%%%%%%%%%%%%%%%%%%%%%%%%%%%%%%%%%%%%%%%%%%%%%%%%%%%%%%%

\begin{theorem}
Every lattice homomorphism is order-preserving.
\end{theorem}

\begin{proof}
Suppose

\[ x\leq y. \]

Then

\[ x=x\wedge y. \]

Applying \(f\),

\[ f(x)=f(x\wedge y)=f(x)\wedge f(y). \]

Therefore

\[ f(x)\leq f(y). \]
\end{proof}

The converse is not generally true: an order-preserving map need not
preserve meets and joins.

%%%%%%%%%%%%%%%%%%%%%%%%%%%%%%%%%%%%%%%%%%%%%%%%%%%%%%%%%%%%
\subsection{Lattice Isomorphisms}
\label{subsec:lattice-isomorphisms}
%%%%%%%%%%%%%%%%%%%%%%%%%%%%%%%%%%%%%%%%%%%%%%%%%%%%%%%%%%%%

A bijective lattice homomorphism is a \emph{lattice isomorphism}.

If \(L\) and \(M\) are isomorphic, write

\[ L\cong M. \]

An isomorphism preserves the complete lattice structure, including order,
meet, and join.

For finite lattices, an isomorphism corresponds to an order-preserving
relabeling of the Hasse diagram.

%%%%%%%%%%%%%%%%%%%%%%%%%%%%%%%%%%%%%%%%%%%%%%%%%%%%%%%%%%%%
\subsection{Congruences on Lattices}
\label{subsec:lattice-congruences}
%%%%%%%%%%%%%%%%%%%%%%%%%%%%%%%%%%%%%%%%%%%%%%%%%%%%%%%%%%%%

Because a lattice is an algebra with operations \(\wedge\) and \(\vee\), a
lattice congruence is an equivalence relation compatible with both.

Thus if

\[ x\mathrel{\theta}x',\qquad y\mathrel{\theta}y', \]

then

\[ x\wedge y\mathrel{\theta}x'\wedge y', \]

and

\[ x\vee y\mathrel{\theta}x'\vee y'. \]

The quotient

\[ L/\theta \]

is again a lattice.

%%%%%%%%%%%%%%%%%%%%%%%%%%%%%%%%%%%%%%%%%%%%%%%%%%%%%%%%%%%%
\subsection{Congruence Lattices}
\label{subsec:lattice-congruence-lattices}
%%%%%%%%%%%%%%%%%%%%%%%%%%%%%%%%%%%%%%%%%%%%%%%%%%%%%%%%%%%%

As introduced in universal algebra, the set of congruences on an algebra
forms a lattice.

For a lattice \(L\), write

\[ \operatorname{Con}(L) \]

for its congruence lattice.

The ordering is inclusion.

The meet of two congruences is their intersection:

\[ \alpha\wedge\beta=\alpha\cap\beta. \]

Their join is the smallest congruence containing both:

\[ \alpha\vee\beta=\operatorname{Cg}(\alpha\cup\beta). \]

%%%%%%%%%%%%%%%%%%%%%%%%%%%%%%%%%%%%%%%%%%%%%%%%%%%%%%%%%%%%
\subsection{A Remarkable Congruence Theorem}
\label{subsec:lattice-congruence-distributive}
%%%%%%%%%%%%%%%%%%%%%%%%%%%%%%%%%%%%%%%%%%%%%%%%%%%%%%%%%%%%

Lattices possess an important universal-algebraic property.

\begin{theorem}
For every lattice \(L\), the congruence lattice
\[ \operatorname{Con}(L) \]
is distributive.
\end{theorem}

Therefore the variety of lattices is \emph{congruence distributive}.

This is a deeper property than saying that the original lattice \(L\) is
distributive.

Indeed, \(L\) itself may be nondistributive while \(\operatorname{Con}(L)\)
is always distributive.

%%%%%%%%%%%%%%%%%%%%%%%%%%%%%%%%%%%%%%%%%%%%%%%%%%%%%%%%%%%%
\subsection{Ideals of a Lattice}
\label{subsec:lattice-ideals}
%%%%%%%%%%%%%%%%%%%%%%%%%%%%%%%%%%%%%%%%%%%%%%%%%%%%%%%%%%%%

The word \emph{ideal} has a lattice-theoretic meaning related to, but
distinct from, ring ideals.

\begin{definitionbox}{Lattice Ideal}
A nonempty subset \(I\subseteq L\) is a \emph{lattice ideal} if:
\begin{enumerate}
    \item whenever \(x\in I\) and \(y\leq x\), then \(y\in I\);
    \item whenever \(x,y\in I\), then \(x\vee y\in I\).
\end{enumerate}
\end{definitionbox}

Thus an ideal is downward closed and closed under finite joins.

%%%%%%%%%%%%%%%%%%%%%%%%%%%%%%%%%%%%%%%%%%%%%%%%%%%%%%%%%%%%
\subsection{Filters}
\label{subsec:lattice-filters}
%%%%%%%%%%%%%%%%%%%%%%%%%%%%%%%%%%%%%%%%%%%%%%%%%%%%%%%%%%%%

The dual concept is a filter.

\begin{definitionbox}{Filter}
A nonempty subset \(F\subseteq L\) is a \emph{filter} if:
\begin{enumerate}
    \item whenever \(x\in F\) and \(x\leq y\), then \(y\in F\);
    \item whenever \(x,y\in F\), then \(x\wedge y\in F\).
\end{enumerate}
\end{definitionbox}

Thus filters are upward closed and closed under finite meets.

The ideal/filter relationship is another manifestation of lattice duality.

%%%%%%%%%%%%%%%%%%%%%%%%%%%%%%%%%%%%%%%%%%%%%%%%%%%%%%%%%%%%
\subsection{Principal Ideals and Filters}
\label{subsec:principal-lattice-ideals}
%%%%%%%%%%%%%%%%%%%%%%%%%%%%%%%%%%%%%%%%%%%%%%%%%%%%%%%%%%%%

For \(a\in L\), the \emph{principal ideal} generated by \(a\) is

\[ \downarrow a=\{x\in L:x\leq a\}. \]

The downward arrow \(\downarrow\) is read ``down-set generated by \(a\)''
or ``principal ideal below \(a\).''

The principal filter generated by \(a\) is

\[ \uparrow a=\{x\in L:a\leq x\}. \]

The upward arrow \(\uparrow\) is read ``up-set generated by \(a\)'' or
``principal filter above \(a\).''

%%%%%%%%%%%%%%%%%%%%%%%%%%%%%%%%%%%%%%%%%%%%%%%%%%%%%%%%%%%%
\subsection{The Lattice of Ideals}
\label{subsec:lattice-of-ideals}
%%%%%%%%%%%%%%%%%%%%%%%%%%%%%%%%%%%%%%%%%%%%%%%%%%%%%%%%%%%%

The collection of ideals of a lattice can itself be ordered by inclusion.

Under suitable definitions, this produces another lattice.

This recurring phenomenon is characteristic of lattice theory:

\begin{conceptbox}
Mathematical structures often generate new lattices consisting of their
substructures, congruences, ideals, closed sets, or other ordered families.
\end{conceptbox}

%%%%%%%%%%%%%%%%%%%%%%%%%%%%%%%%%%%%%%%%%%%%%%%%%%%%%%%%%%%%
\subsection{Closure Operators}
\label{subsec:lattice-closure-operators}
%%%%%%%%%%%%%%%%%%%%%%%%%%%%%%%%%%%%%%%%%%%%%%%%%%%%%%%%%%%%

\begin{definitionbox}{Closure Operator}
A \emph{closure operator} on a poset \(P\) is a function
\[ c:P\to P \]
satisfying:
\[
x\leq c(x),
\qquad
x\leq y\Longrightarrow c(x)\leq c(y),
\qquad
c(c(x))=c(x).
\]
\end{definitionbox}

These properties are called:

\begin{itemize}
    \item extensive;
    \item monotone;
    \item idempotent.
\end{itemize}

Closure operators appear in topology, algebra, geometry, logic, and computer
science.

%%%%%%%%%%%%%%%%%%%%%%%%%%%%%%%%%%%%%%%%%%%%%%%%%%%%%%%%%%%%
\subsection{Closed Elements}
\label{subsec:lattice-closed-elements}
%%%%%%%%%%%%%%%%%%%%%%%%%%%%%%%%%%%%%%%%%%%%%%%%%%%%%%%%%%%%

An element \(x\) is \emph{closed} under \(c\) if

\[ c(x)=x. \]

The set of fixed points of a closure operator often forms a complete
lattice.

Examples include:

\begin{itemize}
    \item closed subsets of a topological space;
    \item linear spans of subsets of a vector space;
    \item subgroups generated by subsets;
    \item algebraic closures in algebraic settings.
\end{itemize}

%%%%%%%%%%%%%%%%%%%%%%%%%%%%%%%%%%%%%%%%%%%%%%%%%%%%%%%%%%%%
\subsection{Monotone Maps}
\label{subsec:lattice-monotone-maps}
%%%%%%%%%%%%%%%%%%%%%%%%%%%%%%%%%%%%%%%%%%%%%%%%%%%%%%%%%%%%

\begin{definitionbox}{Monotone Map}
Let \(P\) and \(Q\) be posets. A function
\[ f:P\to Q \]
is \emph{monotone}, or \emph{order-preserving}, if
\[ x\leq y\Longrightarrow f(x)\leq f(y). \]
\end{definitionbox}

Monotone maps need not preserve meets or joins.

They are therefore more general than lattice homomorphisms.

%%%%%%%%%%%%%%%%%%%%%%%%%%%%%%%%%%%%%%%%%%%%%%%%%%%%%%%%%%%%
\subsection{Fixed Points}
\label{subsec:lattice-fixed-points}
%%%%%%%%%%%%%%%%%%%%%%%%%%%%%%%%%%%%%%%%%%%%%%%%%%%%%%%%%%%%

A fixed point of a function

\[ f:L\to L \]

is an element \(x\in L\) satisfying

\[ f(x)=x. \]

Fixed points are important throughout mathematics.

Lattice theory provides powerful existence theorems when \(L\) is complete
and \(f\) is monotone.

%%%%%%%%%%%%%%%%%%%%%%%%%%%%%%%%%%%%%%%%%%%%%%%%%%%%%%%%%%%%
\subsection{Knaster--Tarski Fixed-Point Theorem}
\label{subsec:knaster-tarski}
%%%%%%%%%%%%%%%%%%%%%%%%%%%%%%%%%%%%%%%%%%%%%%%%%%%%%%%%%%%%

\begin{theorem}[Knaster--Tarski Fixed-Point Theorem]
Let \(L\) be a complete lattice and let
\[ f:L\to L \]
be monotone. Then the set of fixed points of \(f\) forms a complete lattice.
In particular, \(f\) has a least fixed point and a greatest fixed point.
\end{theorem}

This is one of the most important applications of complete lattices.

It has major uses in:

\begin{itemize}
    \item logic;
    \item recursive definitions;
    \item semantics of programming languages;
    \item game theory;
    \item dynamical constructions;
    \item theoretical computer science.
\end{itemize}

%%%%%%%%%%%%%%%%%%%%%%%%%%%%%%%%%%%%%%%%%%%%%%%%%%%%%%%%%%%%
\subsection{Least and Greatest Fixed Points}
\label{subsec:least-greatest-fixed-points}
%%%%%%%%%%%%%%%%%%%%%%%%%%%%%%%%%%%%%%%%%%%%%%%%%%%%%%%%%%%%

The least fixed point can be characterized by

\[ \operatorname{lfp}(f)
=\bigwedge\{x\in L:f(x)\leq x\}. \]

Here

\[ \operatorname{lfp}(f) \]

is read ``the least fixed point of \(f\).''

Similarly,

\[ \operatorname{gfp}(f)
=\bigvee\{x\in L:x\leq f(x)\}. \]

This is read ``the greatest fixed point of \(f\).''

These formulas show how fixed-point theory is built directly from the meet
and join structure of a complete lattice.

%%%%%%%%%%%%%%%%%%%%%%%%%%%%%%%%%%%%%%%%%%%%%%%%%%%%%%%%%%%%
\subsection{Worked Example: A Fixed Point on a Power Set}
\label{subsec:worked-lattice-fixed-point}
%%%%%%%%%%%%%%%%%%%%%%%%%%%%%%%%%%%%%%%%%%%%%%%%%%%%%%%%%%%%

\begin{workedexamplebox}{Fixed Points of a Monotone Map}

Let

\[ X=\{1,2,3\} \]

and define

\[ f:\mathcal P(X)\to\mathcal P(X),\qquad f(A)=A\cup\{1\}. \]

Determine the fixed points.

\begin{tabularx}{\textwidth}{@{}>{\raggedright\arraybackslash}p{0.42\textwidth}X@{}}
A fixed point satisfies \(f(A)=A\). &
\(A\cup\{1\}=A\)\\
This occurs exactly when \(1\in A\). &
\(\{1\}\subseteq A\)\\
List all such subsets. &
\(\{1\},\{1,2\},\{1,3\},\{1,2,3\}\)
\end{tabularx}

\begin{answerbox}
\[ \boxed{\operatorname{Fix}(f)
=\bigl\{\{1\},\{1,2\},\{1,3\},\{1,2,3\}\bigr\}.} \]
\end{answerbox}
\end{workedexamplebox}

%%%%%%%%%%%%%%%%%%%%%%%%%%%%%%%%%%%%%%%%%%%%%%%%%%%%%%%%%%%%
\subsection{Galois Connections}
\label{subsec:lattice-galois-connections}
%%%%%%%%%%%%%%%%%%%%%%%%%%%%%%%%%%%%%%%%%%%%%%%%%%%%%%%%%%%%

The term \emph{Galois connection} refers to an order-theoretic construction
that is broader than Galois theory of field extensions.

\begin{definitionbox}{Galois Connection}
Let \(P\) and \(Q\) be posets. A pair of maps
\[ f:P\to Q,\qquad g:Q\to P \]
forms a \emph{Galois connection} if
\[ f(p)\leq q\iff p\leq g(q). \]
\end{definitionbox}

The two maps determine one another through the order relation.

This idea appears in:

\begin{itemize}
    \item algebra;
    \item topology;
    \item mathematical logic;
    \item formal concept analysis;
    \item category theory.
\end{itemize}

%%%%%%%%%%%%%%%%%%%%%%%%%%%%%%%%%%%%%%%%%%%%%%%%%%%%%%%%%%%%
\subsection{Adjoint Viewpoint}
\label{subsec:lattice-adjoint-viewpoint}
%%%%%%%%%%%%%%%%%%%%%%%%%%%%%%%%%%%%%%%%%%%%%%%%%%%%%%%%%%%%

In a Galois connection,

\[ f(p)\leq q\iff p\leq g(q), \]

the map \(f\) is called a \emph{left adjoint} and \(g\) a \emph{right
adjoint} in the order-theoretic sense.

This is the poset version of the categorical notion of adjunction developed
later in \cref{sec:category-theory}.

A fundamental principle is:

\begin{importantbox}
Left adjoints preserve joins, while right adjoints preserve meets whenever
the relevant joins and meets exist.
\end{importantbox}

%%%%%%%%%%%%%%%%%%%%%%%%%%%%%%%%%%%%%%%%%%%%%%%%%%%%%%%%%%%%
\subsection{Order Reversal and Antitone Maps}
\label{subsec:antitone-maps}
%%%%%%%%%%%%%%%%%%%%%%%%%%%%%%%%%%%%%%%%%%%%%%%%%%%%%%%%%%%%

A function \(f:P\to Q\) is \emph{order-reversing}, or \emph{antitone}, if

\[ x\leq y\Longrightarrow f(y)\leq f(x). \]

Complementation in a Boolean algebra is order-reversing:

\[ x\leq y\Longrightarrow y'\leq x'. \]

This reflects the familiar set-theoretic rule

\[ A\subseteq B\Longrightarrow B^c\subseteq A^c. \]

%%%%%%%%%%%%%%%%%%%%%%%%%%%%%%%%%%%%%%%%%%%%%%%%%%%%%%%%%%%%
\subsection{Atoms and Coatoms}
\label{subsec:atoms-coatoms}
%%%%%%%%%%%%%%%%%%%%%%%%%%%%%%%%%%%%%%%%%%%%%%%%%%%%%%%%%%%%

Let \(L\) be a bounded lattice.

\begin{definitionbox}{Atom}
An \emph{atom} is an element \(a\neq0\) that covers \(0\).
\end{definitionbox}

Thus there is no element strictly between \(0\) and \(a\).

Dually:

\begin{definitionbox}{Coatom}
A \emph{coatom} is an element \(c\neq1\) such that \(1\) covers \(c\).
\end{definitionbox}

In the power-set lattice \(\mathcal P(X)\), the atoms are the singleton
sets.

%%%%%%%%%%%%%%%%%%%%%%%%%%%%%%%%%%%%%%%%%%%%%%%%%%%%%%%%%%%%
\subsection{Atomic Lattices}
\label{subsec:atomic-lattices}
%%%%%%%%%%%%%%%%%%%%%%%%%%%%%%%%%%%%%%%%%%%%%%%%%%%%%%%%%%%%

A bounded lattice is \emph{atomic} if every nonzero element lies above an
atom.

A stronger condition sometimes used is \emph{atomistic}: every element is
a join of atoms.

For a finite set \(X\),

\[ \mathcal P(X) \]

is atomistic because every subset is the union, hence join, of its singleton
elements.

%%%%%%%%%%%%%%%%%%%%%%%%%%%%%%%%%%%%%%%%%%%%%%%%%%%%%%%%%%%%
\subsection{Intervals}
\label{subsec:lattice-intervals}
%%%%%%%%%%%%%%%%%%%%%%%%%%%%%%%%%%%%%%%%%%%%%%%%%%%%%%%%%%%%

If \(a\leq b\) in a lattice \(L\), define the interval

\[ [a,b]=\{x\in L:a\leq x\leq b\}. \]

This is read ``the interval from \(a\) to \(b\).''

Intervals often inherit lattice structure from \(L\).

In particular, if \(x,y\in[a,b]\), then

\[ x\wedge y\in[a,b],\qquad x\vee y\in[a,b]. \]

%%%%%%%%%%%%%%%%%%%%%%%%%%%%%%%%%%%%%%%%%%%%%%%%%%%%%%%%%%%%
\subsection{Finite Lattice Height}
\label{subsec:lattice-height}
%%%%%%%%%%%%%%%%%%%%%%%%%%%%%%%%%%%%%%%%%%%%%%%%%%%%%%%%%%%%

The \emph{height} of a finite lattice is related to the length of its longest
chains.

A chain

\[ x_0<x_1<\cdots<x_n \]

has length \(n\).

Depending on convention, the height of a lattice is the maximum chain length
or the number of levels. The convention should therefore always be stated
when exact numerical height matters.

%%%%%%%%%%%%%%%%%%%%%%%%%%%%%%%%%%%%%%%%%%%%%%%%%%%%%%%%%%%%
\subsection{Ranked Lattices}
\label{subsec:ranked-lattices}
%%%%%%%%%%%%%%%%%%%%%%%%%%%%%%%%%%%%%%%%%%%%%%%%%%%%%%%%%%%%

A finite lattice may admit a rank function

\[ \rho:L\to\N \]

such that covering relations increase rank by one.

The symbol

\[ \rho \]

is the lowercase Greek letter \emph{rho}, pronounced ``row.''

If \(y\) covers \(x\), then

\[ \rho(y)=\rho(x)+1. \]

The Boolean lattice \(B_n\) is naturally ranked by

\[ \rho(A)=|A|. \]

%%%%%%%%%%%%%%%%%%%%%%%%%%%%%%%%%%%%%%%%%%%%%%%%%%%%%%%%%%%%
\subsection{Boolean Lattice Levels}
\label{subsec:boolean-lattice-levels}
%%%%%%%%%%%%%%%%%%%%%%%%%%%%%%%%%%%%%%%%%%%%%%%%%%%%%%%%%%%%

For

\[ B_n=\mathcal P(\{1,\ldots,n\}), \]

the elements at rank \(k\) are precisely the \(k\)-element subsets.

Therefore the number of elements at level \(k\) is

\[ \binom{n}{k}. \]

This connects lattice theory directly with combinatorics.

The levels of \(B_n\) have sizes

\[ \binom n0,\binom n1,\ldots,\binom nn. \]

%%%%%%%%%%%%%%%%%%%%%%%%%%%%%%%%%%%%%%%%%%%%%%%%%%%%%%%%%%%%
\subsection{Lattices and Combinatorics}
\label{subsec:lattices-combinatorics}
%%%%%%%%%%%%%%%%%%%%%%%%%%%%%%%%%%%%%%%%%%%%%%%%%%%%%%%%%%%%

Finite posets and lattices are central objects in combinatorics.

Important questions include:

\begin{itemize}
    \item maximum chain size;
    \item maximum antichain size;
    \item rank enumeration;
    \item Möbius inversion;
    \item incidence algebras;
    \item order ideals;
    \item distributive-lattice representations.
\end{itemize}

These topics are developed more fully in the discrete mathematics and
combinatorics chapters.

%%%%%%%%%%%%%%%%%%%%%%%%%%%%%%%%%%%%%%%%%%%%%%%%%%%%%%%%%%%%
\subsection{Finite Distributive Lattices}
\label{subsec:finite-distributive-lattices}
%%%%%%%%%%%%%%%%%%%%%%%%%%%%%%%%%%%%%%%%%%%%%%%%%%%%%%%%%%%%

Finite distributive lattices possess a particularly elegant representation.

Let \(P\) be a finite poset.

An \emph{order ideal} of \(P\) is a subset \(I\subseteq P\) such that

\[ x\in I,\ y\leq x\Longrightarrow y\in I. \]

Let

\[ J(P) \]

denote the collection of order ideals of \(P\).

Ordered by inclusion, \(J(P)\) is a distributive lattice.

%%%%%%%%%%%%%%%%%%%%%%%%%%%%%%%%%%%%%%%%%%%%%%%%%%%%%%%%%%%%
\subsection{Birkhoff Representation Theorem}
\label{subsec:birkhoff-representation-lattices}
%%%%%%%%%%%%%%%%%%%%%%%%%%%%%%%%%%%%%%%%%%%%%%%%%%%%%%%%%%%%

\begin{theorem}[Birkhoff Representation Theorem]
Every finite distributive lattice is isomorphic to the lattice of order
ideals of a finite poset.
\end{theorem}

More specifically, if \(L\) is finite and distributive, it can be recovered
from the poset of its join-irreducible elements.

This theorem reveals that finite distributive lattices and finite posets are
closely intertwined.

%%%%%%%%%%%%%%%%%%%%%%%%%%%%%%%%%%%%%%%%%%%%%%%%%%%%%%%%%%%%
\subsection{Join-Irreducible Elements}
\label{subsec:join-irreducible}
%%%%%%%%%%%%%%%%%%%%%%%%%%%%%%%%%%%%%%%%%%%%%%%%%%%%%%%%%%%%

\begin{definitionbox}{Join-Irreducible Element}
In a lattice with least element \(0\), a nonzero element \(j\) is
\emph{join-irreducible} if
\[ j=x\vee y\Longrightarrow j=x\text{ or }j=y. \]
\end{definitionbox}

Such an element cannot be decomposed nontrivially as a join of smaller
elements.

Dually, a \emph{meet-irreducible} element cannot be decomposed nontrivially
as a meet.

%%%%%%%%%%%%%%%%%%%%%%%%%%%%%%%%%%%%%%%%%%%%%%%%%%%%%%%%%%%%
\subsection{Lattice Theory and Topology}
\label{subsec:lattice-topology}
%%%%%%%%%%%%%%%%%%%%%%%%%%%%%%%%%%%%%%%%%%%%%%%%%%%%%%%%%%%%

For a topological space \(X\), the collection of open sets

\[ \mathcal O(X) \]

is ordered by inclusion.

Arbitrary unions of open sets are open, and finite intersections are open.

Thus

\[ \bigvee_{i\in I}U_i=\bigcup_{i\in I}U_i, \]

while finite meets are

\[ U\wedge V=U\cap V. \]

The lattice of open sets is an example of a more specialized structure
called a \emph{frame}.

This connects lattice theory with point-set topology and locale theory.

%%%%%%%%%%%%%%%%%%%%%%%%%%%%%%%%%%%%%%%%%%%%%%%%%%%%%%%%%%%%
\subsection{Frames}
\label{subsec:lattice-frames}
%%%%%%%%%%%%%%%%%%%%%%%%%%%%%%%%%%%%%%%%%%%%%%%%%%%%%%%%%%%%

\begin{definitionbox}{Frame}
A \emph{frame} is a complete lattice in which finite meets distribute over
arbitrary joins:
\[ x\wedge\left(\bigvee_{i\in I}y_i\right)
=\bigvee_{i\in I}(x\wedge y_i). \]
\end{definitionbox}

The lattice of open sets of every topological space is a frame.

Frames allow aspects of topology to be studied algebraically.

%%%%%%%%%%%%%%%%%%%%%%%%%%%%%%%%%%%%%%%%%%%%%%%%%%%%%%%%%%%%
\subsection{Lattice Theory and Logic}
\label{subsec:lattice-logic}
%%%%%%%%%%%%%%%%%%%%%%%%%%%%%%%%%%%%%%%%%%%%%%%%%%%%%%%%%%%%

Logical propositions can be ordered by implication.

If

\[ P\Longrightarrow Q, \]

one may regard \(P\) as lying below \(Q\).

Then conjunction behaves like meet:

\[ P\wedge Q, \]

and disjunction behaves like join:

\[ P\vee Q. \]

Boolean algebras model classical propositional logic.

Other lattice-like structures model intuitionistic and nonclassical logics.

%%%%%%%%%%%%%%%%%%%%%%%%%%%%%%%%%%%%%%%%%%%%%%%%%%%%%%%%%%%%
\subsection{Lattice Theory and Algebra}
\label{subsec:lattice-algebra}
%%%%%%%%%%%%%%%%%%%%%%%%%%%%%%%%%%%%%%%%%%%%%%%%%%%%%%%%%%%%

Many algebraic structures naturally produce lattices.

\begin{center}
\begin{tabularx}{0.92\textwidth}{@{}lXX@{}}
\toprule
\textbf{Structure} & \textbf{Associated Lattice} & \textbf{Order}\\
\midrule
Group \(G\) & subgroups of \(G\) & inclusion\\
Ring \(R\) & ideals of \(R\) & inclusion\\
Module \(M\) & submodules of \(M\) & inclusion\\
Vector space \(V\) & subspaces of \(V\) & inclusion\\
Algebra \(\mathbf A\) & congruences of \(\mathbf A\) & inclusion\\
Set \(X\) & subsets of \(X\) & inclusion\\
\bottomrule
\end{tabularx}
\end{center}

Lattice theory therefore provides a common language for studying systems of
subobjects.

%%%%%%%%%%%%%%%%%%%%%%%%%%%%%%%%%%%%%%%%%%%%%%%%%%%%%%%%%%%%
\subsection{Lattice Theory and Data}
\label{subsec:lattice-data}
%%%%%%%%%%%%%%%%%%%%%%%%%%%%%%%%%%%%%%%%%%%%%%%%%%%%%%%%%%%%

Order structures also appear in data analysis and knowledge representation.

For example, formal concept analysis organizes relationships between
objects and attributes into a \emph{concept lattice}.

Concept lattices reveal hierarchical patterns and dependencies within data.

This is one example of lattice theory extending beyond traditional pure
algebra.

%%%%%%%%%%%%%%%%%%%%%%%%%%%%%%%%%%%%%%%%%%%%%%%%%%%%%%%%%%%%
\subsection{A Hierarchy of Lattice Structures}
\label{subsec:lattice-hierarchy}
%%%%%%%%%%%%%%%%%%%%%%%%%%%%%%%%%%%%%%%%%%%%%%%%%%%%%%%%%%%%

Several important classes fit into the hierarchy

\begin{center}
\begin{tikzpicture}[
    node distance=8mm,
    box/.style={draw,rounded corners,align=center,inner sep=5pt},
    >=stealth
]
\node[box] (poset) {Posets};
\node[box,below=of poset] (lattice) {Lattices};
\node[box,below=of lattice] (bounded) {Bounded lattices};
\node[box,below left=of bounded] (modular) {Modular lattices};
\node[box,below right=of bounded] (comp) {Complemented lattices};
\node[box,below=18mm of bounded] (dist) {Distributive lattices};
\node[box,below=of dist] (bool) {Boolean algebras};

\draw[->] (lattice)--(poset);
\draw[->] (bounded)--(lattice);
\draw[->] (modular)--(lattice);
\draw[->] (dist)--(modular);
\draw[->] (comp)--(bounded);
\draw[->] (bool)--(dist);
\draw[->] (bool)--(comp);
\end{tikzpicture}
\end{center}

The diagram indicates implications, not a complete classification of every
possible relationship.

%%%%%%%%%%%%%%%%%%%%%%%%%%%%%%%%%%%%%%%%%%%%%%%%%%%%%%%%%%%%
\subsection{Common Errors}
\label{subsec:lattice-common-errors}
%%%%%%%%%%%%%%%%%%%%%%%%%%%%%%%%%%%%%%%%%%%%%%%%%%%%%%%%%%%%

\subsubsection{Assuming Every Poset Is a Lattice}

A lattice requires every pair to possess both a meet and a join.

A general poset may lack either.

%%%%%%%%%%%%%%%%%%%%%%%%%%%%%%%%%%%%%%%%%%%%%%%%%%%%%%%%%%%%
\subsubsection{Confusing Minimal with Least}

A minimal element has nothing strictly below it.

A least element lies below every element.

There may be many minimal elements but at most one least element.

%%%%%%%%%%%%%%%%%%%%%%%%%%%%%%%%%%%%%%%%%%%%%%%%%%%%%%%%%%%%
\subsubsection{Confusing Maximal with Greatest}

A maximal element has nothing strictly above it.

A greatest element lies above every element.

These notions are not equivalent.

%%%%%%%%%%%%%%%%%%%%%%%%%%%%%%%%%%%%%%%%%%%%%%%%%%%%%%%%%%%%
\subsubsection{Assuming Meet Means Minimum}

The meet

\[ x\wedge y \]

is the greatest lower bound of \(x\) and \(y\).

If \(x\) and \(y\) are incomparable, the meet may be different from both.

%%%%%%%%%%%%%%%%%%%%%%%%%%%%%%%%%%%%%%%%%%%%%%%%%%%%%%%%%%%%
\subsubsection{Assuming Join Means Maximum}

The join

\[ x\vee y \]

is the least upper bound.

It need not equal \(x\) or \(y\).

%%%%%%%%%%%%%%%%%%%%%%%%%%%%%%%%%%%%%%%%%%%%%%%%%%%%%%%%%%%%
\subsubsection{Reading a Hasse Diagram as an Ordinary Graph}

Edges in a Hasse diagram represent covering relations.

Transitive order relations are implied even when no edge is drawn.

%%%%%%%%%%%%%%%%%%%%%%%%%%%%%%%%%%%%%%%%%%%%%%%%%%%%%%%%%%%%
\subsubsection{Including Transitive Edges in a Hasse Diagram}

If

\[ x<y<z, \]

the relation \(x<z\) is already implied.

A Hasse diagram normally does not draw a direct edge from \(x\) to \(z\).

%%%%%%%%%%%%%%%%%%%%%%%%%%%%%%%%%%%%%%%%%%%%%%%%%%%%%%%%%%%%
\subsubsection{Assuming Every Bounded Lattice Is Complemented}

Having \(0\) and \(1\) does not guarantee that every element has a
complement.

%%%%%%%%%%%%%%%%%%%%%%%%%%%%%%%%%%%%%%%%%%%%%%%%%%%%%%%%%%%%
\subsubsection{Assuming Complements Are Always Unique}

Complements need not be unique in arbitrary complemented lattices.

They are unique in Boolean algebras because distributivity is available.

%%%%%%%%%%%%%%%%%%%%%%%%%%%%%%%%%%%%%%%%%%%%%%%%%%%%%%%%%%%%
\subsubsection{Assuming Modular Means Distributive}

Every distributive lattice is modular, but not every modular lattice is
distributive.

The lattice \(M_3\) is the standard counterexample.

%%%%%%%%%%%%%%%%%%%%%%%%%%%%%%%%%%%%%%%%%%%%%%%%%%%%%%%%%%%%
\subsubsection{Assuming a Lattice Is Automatically Complete}

A lattice only guarantees binary meets and joins.

A complete lattice requires meets and joins for every subset.

%%%%%%%%%%%%%%%%%%%%%%%%%%%%%%%%%%%%%%%%%%%%%%%%%%%%%%%%%%%%
\subsubsection{Confusing Lattice Ideals with Ring Ideals}

The terminology is related historically and structurally, but the
definitions are different.

A lattice ideal is a downward-closed subset closed under joins.

%%%%%%%%%%%%%%%%%%%%%%%%%%%%%%%%%%%%%%%%%%%%%%%%%%%%%%%%%%%%
\subsubsection{Assuming Every Monotone Map Is a Lattice Homomorphism}

A monotone map preserves order.

A lattice homomorphism must preserve both meet and join.

The latter is a stronger requirement.

%%%%%%%%%%%%%%%%%%%%%%%%%%%%%%%%%%%%%%%%%%%%%%%%%%%%%%%%%%%%
\subsection{Applications}
\label{subsec:lattice-applications}
%%%%%%%%%%%%%%%%%%%%%%%%%%%%%%%%%%%%%%%%%%%%%%%%%%%%%%%%%%%%

\begin{applicationbox}
\textbf{Abstract Algebra.}
Subgroups, ideals, submodules, subspaces, and congruences naturally form
lattices under inclusion.

\medskip

\textbf{Universal Algebra.}
Congruence lattices encode how an algebra can be partitioned into compatible
quotient structures.

\medskip

\textbf{Logic.}
Meet, join, and complement model conjunction, disjunction, and negation.
Boolean algebras provide an algebraic model of classical propositional
logic.

\medskip

\textbf{Set Theory.}
Power sets are complete Boolean lattices under inclusion, intersection,
union, and complement.

\medskip

\textbf{Linear Algebra.}
Subspaces form modular lattices, with intersection as meet and subspace sum
as join.

\medskip

\textbf{Topology.}
Open sets form complete lattice-like structures called frames, linking
lattice theory with point-free topology.

\medskip

\textbf{Combinatorics.}
Finite lattices interact with chains, antichains, rank functions, Boolean
lattices, Möbius inversion, and incidence structures.

\medskip

\textbf{Computer Science.}
Complete lattices and fixed-point theorems are used in program semantics,
static analysis, recursive definitions, databases, and formal verification.

\medskip

\textbf{Data Analysis.}
Formal concept analysis uses concept lattices to organize relationships
between objects and attributes.
\end{applicationbox}

%%%%%%%%%%%%%%%%%%%%%%%%%%%%%%%%%%%%%%%%%%%%%%%%%%%%%%%%%%%%
\subsection{Exercises}
\label{subsec:lattice-exercises}
%%%%%%%%%%%%%%%%%%%%%%%%%%%%%%%%%%%%%%%%%%%%%%%%%%%%%%%%%%%%

\begin{exercise}[1]
Define a partially ordered set.
\end{exercise}

\begin{exercise}[1]
What does it mean for two elements of a poset to be comparable?
\end{exercise}

\begin{exercise}[1]
Define a chain.
\end{exercise}

\begin{exercise}[1]
Define an antichain.
\end{exercise}

\begin{exercise}[1]
What does it mean for \(y\) to cover \(x\)?
\end{exercise}

\begin{exercise}[1]
Explain which relations are omitted from a Hasse diagram.
\end{exercise}

\begin{exercise}[2]
Draw the Hasse diagram of the divisors of \(18\) ordered by divisibility.
\end{exercise}

\begin{exercise}[2]
Draw the Hasse diagram of the divisors of \(24\) ordered by divisibility.
\end{exercise}

\begin{exercise}[1]
Define a minimal element.
\end{exercise}

\begin{exercise}[1]
Define a least element.
\end{exercise}

\begin{exercise}[2]
Give an example of a poset having more than one minimal element but no least
element.
\end{exercise}

\begin{exercise}[1]
Define a maximal element.
\end{exercise}

\begin{exercise}[1]
Define a greatest element.
\end{exercise}

\begin{exercise}[1]
Define an upper bound of a subset.
\end{exercise}

\begin{exercise}[1]
Define a lower bound of a subset.
\end{exercise}

\begin{exercise}[1]
Define supremum and infimum.
\end{exercise}

\begin{exercise}[1]
Define the join \(x\vee y\).
\end{exercise}

\begin{exercise}[1]
Define the meet \(x\wedge y\).
\end{exercise}

\begin{exercise}[2]
In the divisibility lattice of the divisors of \(48\), compute
\[ 8\wedge12,\qquad8\vee12. \]
\end{exercise}

\begin{exercise}[2]
In the divisibility lattice of the divisors of \(72\), compute
\[ 9\wedge12,\qquad9\vee12. \]
\end{exercise}

\begin{exercise}[1]
State the order-theoretic definition of a lattice.
\end{exercise}

\begin{exercise}[2]
Give an example of a finite poset that is not a lattice and explain which
meet or join fails to exist.
\end{exercise}

\begin{exercise}[2]
Show that
\[ \mathcal P(X) \]
is a lattice under inclusion.
\end{exercise}

\begin{exercise}[2]
Show that in \(\mathcal P(X)\),
\[ A\wedge B=A\cap B,\qquad A\vee B=A\cup B. \]
\end{exercise}

\begin{exercise}[2]
Show that the divisors of a fixed positive integer \(n\) form a lattice
under divisibility.
\end{exercise}

\begin{exercise}[2]
For subspaces \(U,W\leq V\), explain why
\[ U\cap W \]
is the meet and
\[ U+W \]
is the join.
\end{exercise}

\begin{exercise}[2]
For subgroups \(H,K\leq G\), explain why the join is generally
\[ \langle H\cup K\rangle \]
rather than \(H\cup K\).
\end{exercise}

\begin{exercise}[1]
State the algebraic lattice identities.
\end{exercise}

\begin{exercise}[2]
Verify the absorption laws for sets using union and intersection.
\end{exercise}

\begin{exercise}[2]
Prove that
\[ x\leq y\iff x\wedge y=x. \]
\end{exercise}

\begin{exercise}[2]
Prove that
\[ x\leq y\iff x\vee y=y. \]
\end{exercise}

\begin{exercise}[1]
State the principle of lattice duality.
\end{exercise}

\begin{exercise}[2]
Write the dual of
\[ x\wedge0=0. \]
\end{exercise}

\begin{exercise}[1]
Define a sublattice.
\end{exercise}

\begin{exercise}[2]
Give an example of a subset of a lattice that is a subposet but not a
sublattice.
\end{exercise}

\begin{exercise}[1]
Define a bounded lattice.
\end{exercise}

\begin{exercise}[1]
Define a complement of an element.
\end{exercise}

\begin{exercise}[1]
Define a complemented lattice.
\end{exercise}

\begin{exercise}[2]
Find the complement of each element of
\[ \mathcal P(\{1,2,3\}). \]
\end{exercise}

\begin{exercise}[1]
Define a distributive lattice.
\end{exercise}

\begin{exercise}[2]
Verify that the power-set lattice is distributive.
\end{exercise}

\begin{exercise}[2]
Use \(M_3\) to demonstrate explicitly that distributivity can fail.
\end{exercise}

\begin{exercise}[2]
Explain why \(N_5\) is called the pentagon lattice.
\end{exercise}

\begin{exercise}[2]
State Birkhoff's characterization of distributive lattices using \(M_3\) and
\(N_5\).
\end{exercise}

\begin{exercise}[1]
Define a modular lattice.
\end{exercise}

\begin{exercise}[2]
Explain why every distributive lattice is modular.
\end{exercise}

\begin{exercise}[2]
Give an example of a modular lattice that is not distributive.
\end{exercise}

\begin{exercise}[2]
State the characterization of modular lattices using \(N_5\).
\end{exercise}

\begin{exercise}[2]
Verify the modular law for subspaces of a vector space in a simple example.
\end{exercise}

\begin{exercise}[1]
Define a Boolean algebra.
\end{exercise}

\begin{exercise}[2]
Explain why
\[ \mathcal P(X) \]
is a Boolean algebra.
\end{exercise}

\begin{exercise}[2]
Prove one of De Morgan's laws in a Boolean algebra.
\end{exercise}

\begin{exercise}[3]
Prove that complements are unique in a distributive complemented lattice.
\end{exercise}

\begin{exercise}[1]
Define a complete lattice.
\end{exercise}

\begin{exercise}[2]
Explain why every complete lattice is bounded.
\end{exercise}

\begin{exercise}[2]
Show that
\[ \mathcal P(X) \]
is a complete lattice.
\end{exercise}

\begin{exercise}[2]
For a family \(\mathcal A\subseteq\mathcal P(X)\), identify
\[ \bigvee\mathcal A\qquad\text{and}\qquad\bigwedge\mathcal A. \]
\end{exercise}

\begin{exercise}[1]
Define a lattice homomorphism.
\end{exercise}

\begin{exercise}[2]
Prove that every lattice homomorphism is monotone.
\end{exercise}

\begin{exercise}[2]
Give an example of a monotone map that is not a lattice homomorphism.
\end{exercise}

\begin{exercise}[1]
Define a lattice congruence.
\end{exercise}

\begin{exercise}[2]
Explain why a quotient by a lattice congruence is again a lattice.
\end{exercise}

\begin{exercise}[1]
Define a lattice ideal.
\end{exercise}

\begin{exercise}[1]
Define a lattice filter.
\end{exercise}

\begin{exercise}[2]
Explain how ideals and filters are dual notions.
\end{exercise}

\begin{exercise}[1]
Define the principal ideal
\[ \downarrow a. \]
\end{exercise}

\begin{exercise}[1]
Define the principal filter
\[ \uparrow a. \]
\end{exercise}

\begin{exercise}[1]
Define a closure operator.
\end{exercise}

\begin{exercise}[2]
Show that the map
\[ A\mapsto\Span(A) \]
acts as a closure operator on subsets of a vector space.
\end{exercise}

\begin{exercise}[1]
Define a monotone map.
\end{exercise}

\begin{exercise}[1]
Define a fixed point.
\end{exercise}

\begin{exercise}[2]
State the Knaster--Tarski fixed-point theorem.
\end{exercise}

\begin{exercise}[2]
Let
\[ f(A)=A\cup\{1,2\} \]
on \(\mathcal P(\{1,2,3,4\})\). Determine all fixed points.
\end{exercise}

\begin{exercise}[2]
For the previous function, determine the least and greatest fixed points.
\end{exercise}

\begin{exercise}[1]
Define a Galois connection.
\end{exercise}

\begin{exercise}[2]
Explain what it means for one map to be a left adjoint and another to be a
right adjoint between posets.
\end{exercise}

\begin{exercise}[2]
State which structures are preserved by left and right adjoints.
\end{exercise}

\begin{exercise}[1]
Define an atom of a bounded lattice.
\end{exercise}

\begin{exercise}[1]
Define a coatom.
\end{exercise}

\begin{exercise}[2]
Identify the atoms and coatoms of
\[ \mathcal P(\{1,2,3,4\}). \]
\end{exercise}

\begin{exercise}[1]
Define an interval \([a,b]\) in a lattice.
\end{exercise}

\begin{exercise}[2]
Show that an interval in a lattice is itself a sublattice.
\end{exercise}

\begin{exercise}[1]
What is a ranked lattice?
\end{exercise}

\begin{exercise}[2]
Show that the Boolean lattice \(B_n\) is ranked by
\[ \rho(A)=|A|. \]
\end{exercise}

\begin{exercise}[2]
How many elements occur at rank \(k\) in \(B_n\)?
\end{exercise}

\begin{exercise}[1]
Define a join-irreducible element.
\end{exercise}

\begin{exercise}[2]
Find the join-irreducible elements of \(B_3\).
\end{exercise}

\begin{exercise}[2]
State the finite Birkhoff representation theorem.
\end{exercise}

\begin{exercise}[1]
Define a frame.
\end{exercise}

\begin{exercise}[2]
Explain why the open subsets of a topological space form a frame.
\end{exercise}

\begin{exercise}[2]
Explain how Boolean algebra models classical propositional logic.
\end{exercise}

\begin{exercise}[2]
Explain how subgroup lattices connect group theory with lattice theory.
\end{exercise}

\begin{exercise}[2]
Explain why congruence lattices are important in universal algebra.
\end{exercise}

\begin{exercise}[3]
Compare the following structures:
\[
\text{bounded},\qquad
\text{complemented},\qquad
\text{modular},\qquad
\text{distributive},\qquad
\text{Boolean}.
\]
State which implications always hold and which do not.
\end{exercise}

%%%%%%%%%%%%%%%%%%%%%%%%%%%%%%%%%%%%%%%%%%%%%%%%%%%%%%%%%%%%
\subsection{Section Connections}
\label{subsec:lattice-theory-connections}
%%%%%%%%%%%%%%%%%%%%%%%%%%%%%%%%%%%%%%%%%%%%%%%%%%%%%%%%%%%%

\chapterconnection{
Lattice theory develops the mathematics of ordered structures in which
pairs of elements possess canonical greatest lower bounds and least upper
bounds. The order-theoretic operations of meet and join admit an equivalent
algebraic description through identities, placing lattices naturally within
universal algebra. Power sets, divisor systems, subgroup systems, ideals,
submodules, subspaces, and congruences all produce important lattice
examples. Distributive, modular, complemented, Boolean, and complete
lattices provide increasingly specialized structures with applications
across algebra, logic, topology, combinatorics, and computer science.
Congruence lattices connect lattice theory directly with universal algebra,
while Galois connections and fixed-point theorems anticipate categorical
adjunctions and order-theoretic methods used in logic and computation.
}

%%%%%%%%%%%%%%%%%%%%%%%%%%%%%%%%%%%%%%%%%%%%%%%%%%%%%%%%%%%%
% SECTION SOURCE TRACKING
%%%%%%%%%%%%%%%%%%%%%%%%%%%%%%%%%%%%%%%%%%%%%%%%%%%%%%%%%%%%

%------------------------------------------------------------
% 1.15 Lattice Theory
%------------------------------------------------------------
%
% Sources:
%   - B. A. Davey and H. A. Priestley,
%     Introduction to Lattices and Order.
%     Key: daveypriestley2002lattices
%
%   - Stanley Burris and H. P. Sankappanavar,
%     A Course in Universal Algebra.
%     Key: burrissankappanavar1981universal
%
% Used for:
%   - posets, chains, antichains, and Hasse diagrams;
%   - meets, joins, bounds, suprema, and infima;
%   - order-theoretic and algebraic lattice definitions;
%   - duality;
%   - sublattices and product lattices;
%   - bounded and complemented lattices;
%   - distributive and modular lattices;
%   - M_3 and N_5;
%   - Boolean algebras;
%   - complete lattices;
%   - lattice homomorphisms;
%   - ideals and filters;
%   - closure operators;
%   - fixed-point theory;
%   - Galois connections;
%   - atoms and ranked lattices;
%   - finite distributive lattice representation;
%   - congruence-lattice connections with universal algebra.
%
% Notes:
%   - Basic relation and partial-order definitions are introduced
%     canonically in Foundations.
%   - Group-specific subgroup lattices belong canonically to Group Theory.
%   - Ring ideal theory belongs canonically to Ring Theory.
%   - Module submodule theory belongs canonically to Module Theory.
%   - Congruence theory is introduced canonically in Universal Algebra.
%   - Detailed order-theoretic combinatorics belongs to Discrete Mathematics
%     and Combinatorics.
%   - Topological frames are introduced only as a connection; topology is
%     developed canonically in Chapter 4.
%   - Categorical adjunctions are developed canonically in Category Theory.
%   - Visualizations and pedagogical organization are independently
%     developed for this text.
%
\nocite{daveypriestley2002lattices}
\nocite{burrissankappanavar1981universal}